\section{引言}
\label{sec:intro}

无线自组织网络在物联网、车载通信与无人机集群等场景中已成为关键的承载技术，
其 MAC 层的调度策略直接决定网络的吞吐量、时延与节点公平性~\cite{author2024survey}。
时分多址（TDMA）通过将信道划分为固定帧和时隙、为节点分配无冲突的接入机会，
是上述场景中实现确定性低时延的主流方案，并可通过空间复用进一步提高频谱效率~\cite{spatial_reuse}。
然而，节点移动、链路开断与业务异构使可调度时隙的分布持续漂移，
原本针对静态拓扑设计的启发式调度难以同时维持高吞吐与可控的长尾时延。

传统 TDMA 调度算法可分为两类：
\emph{所有者时隙}类方案（plain TDMA 及其增强版本）按固定划分保证无冲突，
但稀疏拓扑下大量时隙闲置而密集区域产生饥饿；
\emph{两跳仲裁与混合分配}类方案（如 Z-MAC~\cite{zmac}、TRAMA~\cite{trama}、STDMA~\cite{stdma}）
通过混合分配/竞争或确定性两跳仲裁缓解上述问题，
但其权值更新仍依赖本地参数手工调优，
在到达过程剧烈波动与拓扑突变下难以兼顾吞吐、公平性与扰动恢复速度。
近年来兴起的基于深度强化学习的 MAC 调度~\cite{drlmac2019,ppoma2023,mnih2015dqn}
为动态调度提供了数据驱动的替代路径。
然而，本文在沿用这一思路构造朴素基线时观察到：
若直接由 PPO 端到端输出申请概率（本文将该朴素基线记为 \textbf{PPO-Base}），
会带来两类新风险：
其一，学习型策略容易破坏仿真器中已有启发式调度器的安全退避先验，
在训练初期反而恶化吞吐与公平性；
其二，对多维 Bernoulli 动作仅给出聚合奖励，无法精确归因每个时隙决策的好坏，
导致策略熵塌缩或学习停滞。
本文正是从修复 PPO-Base 的上述两点不足出发，
逐步引出后述的 DASD-PPO 设计。

本文针对上述问题，提出\textbf{密度感知服务债务 PPO}（DASD-PPO）调度框架。
DASD-PPO 在策略空间中显式引入与拓扑感知一致的归纳偏置，
其核心是\emph{将 PPO 学习信号约束在启发式可行域内，并对长尾未服务节点提供确定性补偿}。
具体而言，DASD-PPO 由四个相互配合的机制组成：
（i）\textbf{乘数模式策略合成与启发式偏离锚定}，
让 Actor 输出乘数因子 \(\alpha\in(0,1)^M\)，
通过 \(p_\text{req}=\text{clamp}(p_\text{heur}\cdot 2\alpha,0,1)\) 合成申请概率，
并以与启发式概率的 \(\ell_2\) 距离作为软惩罚锚定策略；
（ii）\textbf{两跳安全掩膜}，
在采样阶段显式屏蔽上一帧一/二跳邻居占用的时隙，
保留拓扑感知的可行域；
（iii）\textbf{服务债务与请求预算耦合}，
对超过门限连续未发送且队列非空的节点施加确定性动作推力，
并以预算上限防止过申请；
（iv）\textbf{密度感知自适应控制}，
按归一化网络密度 \(\rho_t\) 在掩膜启用/关闭之间平滑插值，
使框架能在稀疏与密集拓扑间统一适配。

为进一步定位上述机制的边际贡献，本文构造了从弱到强的三层学习型算法演进：
\textbf{PPO-Base}（基础逐时隙 Bernoulli 采样）、
\textbf{HC-PPO}（启发式约束 PPO，在 PPO-Base 之上加入启发式偏离软惩罚）、
\textbf{DASD-PPO}（在 HC-PPO 之上再叠加两跳掩膜、服务债务、请求预算与密度自适应）。
三者共用相同的 PPO 训练超参数与训练预算，
其差异完全由前述四个机制贡献。
此外，本文实现了 OMNeT++ 仿真器与外部 PPO 训练器的异步在线训练管道：
仿真器与训练器通过两个 POSIX 命名管道完成状态/动作交换，
任何一端未就绪时另一端自动降级或重连，
为在动态拓扑上做长时序在线学习提供工程基础。

本文的主要贡献可概括为：
\begin{itemize}
    \item \textbf{策略框架}：提出 DASD-PPO 调度算法，
          以乘数模式 + 两跳掩膜 + 服务债务 + 密度自适应四项机制为骨架，
          显式约束 RL 策略不偏离拓扑感知可行域，
          并对长尾节点提供确定性补偿。
    \item \textbf{算法机理}：基于 LSTM Actor--Critic（每节点独立、约 13.8K 参数）
          实现轻量级在线策略；针对多维 Bernoulli 动作引入\emph{逐时隙}信用分配，
          使 Critic 输出 \(M\) 维逐时隙价值，PPO 采用逐时隙比率 \(\times\) 优势函数更新，
          显著提升熵收敛速度与样本效率。
    \item \textbf{系统实现}：构建 OMNeT++ 6.3 仿真器与外部 PPO 训练器的
          POSIX 命名管道双向异步交互管道，支持在线训练与启发式降级，
          能在两种 \(N=12\) 主场景与四种动态扰动场景上稳定运行。
    \item \textbf{完备评估}：在两个静态主场景与四个动态扰动场景上做了 3 个随机种子下的
          十种算法横向对比；结果表明 DASD-PPO 在 UAV 稀疏拓扑下取得\emph{全表最低}的饥饿比例
          （\(0.2500\) vs. HC-PPO 的 \(0.5556\)），
          在 Bridge Break 与 Node Rejoin 扰动下分别取得最小 PDR 衰减与最高 throughput 恢复，
          学习型方法整体相对论文启发式近似（Z-MAC*、TRAMA*、Bandit-Index*）
          PDR 与 throughput 提升数倍。
\end{itemize}

本文其余部分组织如下：
\Cref{sec:related} 回顾相关工作；
\Cref{sec:model} 给出系统模型与优化目标的形式化描述；
\Cref{sec:method} 详细介绍 DASD-PPO 的四个机制与训练流程；
\Cref{sec:exp} 报告静态主场景、动态扰动场景及算法演进消融的实验结果；
\Cref{sec:discussion} 讨论设计权衡与局限；
\Cref{sec:conclusion} 总结全文并展望未来工作。
